\documentclass[b5paper,10pt,twoside]{book}
\usepackage[T1]{fontenc}
\usepackage[utf8]{inputenc}
\usepackage[spanish]{babel}
\usepackage[super]{nth}
\usepackage{
    amsmath,
    amsfonts,
    amssymb,
    amsthm,
    imakeidx,
    graphicx,
    polynom,
    % subcaption,
    lipsum,
    lmodern,
    systeme,
    float,
    ifthen,
    mathrsfs,
    units,
    % enumerate,
    enumitem,
    multicol,
    geometry,
    mdframed,
    setspace,
    subcaption,
    anyfontsize,
    mathtools,
    fancyhdr,
    booktabs,
    multirow,
    csquotes,
    microtype,
    bookmark,
    emptypage,
    blkarray,
    aliascnt,
    dsfont,
    comment,
    tabularx,
    % mathds,
    titlesec,
    cancel,
    currfile,
    xcolor,
}
\usepackage[perpage]{footmisc}

\setenumerate{topsep=0pt}

\setlength{\parskip}{0.5\baselineskip}%

% \fontseries{b}\selectfont
% \renewcommand{\bfdefault}{b}

\titleformat{\chapter}
  {\normalfont\Huge}{\Huge\thechapter}{15pt}{\Huge}

% \titleformat{\section}
%   {\normalfont\large}{\large\thechapter}{15pt}{\large}

% \titleformat{\section}
%   {\normalfont\Huge}{\Huge\thechapter}{15pt}{\Huge}

% \onehalfspacing
\setstretch{1.05}
\geometry{
    margin=1in,
    % headheight=50pt,
    % headsep=12pt,
    % showframe
    % textwidth=5in,
}
% \decimalpoint
\raggedbottom
\pagestyle{plain}
% \fancyhf{}
% % \fancyheadoffset{1.0mm}
% \lhead{\includegraphics[width=0.35\textwidth]{logoECMC.png}}
% \rhead{\includegraphics[width=0.25\textwidth]{logoYachay.png}}
\fancyfoot[RO,LE]{\thepage}

\usepackage[delims=\lbrack\rbrack]{spalign}
\let\matrix=\spalignmat
\let\Amatrix=\spalignaugmat
\let\AAmatrix=\spalignaugmatn
\newcommand{\Dmatrix}[1]{\spaligndelims\vert\vert\spalignmat{#1}}

% \fancyhead[RE]{\scshape{\nouppercase{\leftmark}}}
% \fancyhead[LO]{\scshape{\nouppercase{\rightmark}}}
\renewcommand{\headrulewidth}{0pt}
\graphicspath{{images/}}
\bookmarksetup{numbered}
\makeatletter
 \renewcommand\Hy@numberline[1]{#1. }
\makeatother
\usepackage{hyperref}
\hypersetup{%
    colorlinks = {true},
            linkcolor = {blue},
            urlcolor  = {blue},
            citecolor = {blue},
            anchorcolor = {blue},
    pdftitle    ={},
    pdfauthor   ={K. Christian Chávez},
    pdfsubject  ={Abstract Algebra},
    pdfcreator  ={LaTeX},
    pdfproducer ={},
    % pdfkeywords={}{}{},
    % hidelinks,
    pdfpagemode={UseNone},
}
\newcolumntype{Y}{>{\centering\arraybackslash}X}
\renewcommand\tabularxcolumn[1]{m{#1}}
% \renewcommand\arraystretch{1.5}

\newcommand{\p}{\cdot}
\newenvironment{solution}{\noindent\textbf{\textsl{Solution}.}}{\hfill$\square$}
% \newenvironment{remark}{\noindent\textbf{\textsl{Remark.}}}{}
% \theoremstyle{definition}
% \newtheorem{problem}{Problem}%[chapter]

% \setlength\parindent{0pt}

% Nuevos comandos
\let\originalleft\left
\let\originalright\right
% \renewcommand{\left}{\mathopen{}\mathclose\bgroup\originalleft}
% \renewcommand{\right}{\aftergroup\egroup\originalright}
\newcommand{\N}{\mathds{N}}
\newcommand{\Z}{\mathds{Z}}
\newcommand{\R}{\mathds{R}}
\newcommand{\Q}{\mathds{Q}}
\newcommand{\II}{\mathds{I}}
\newcommand{\K}{\mathds{K}}
\newcommand{\PP}{\mathds{P}}
\newcommand{\C}{\mathds{C}}
\newcommand{\T}{\mathscr{T}}
\newcommand{\I}{\mathcal{I}}
\newcommand{\E}{\mathscr{E}}
\newcommand{\B}{\mathcal{B}}
\newcommand{\F}{\mathscr{F}}
\newcommand{\U}{\mathscr{U}}
\newcommand{\LL}{\mathscr{L}}
\newcommand{\OO}{\mathscr{O}}
\newcommand{\Fam}{\mathscr{F}}
\newcommand{\Nei}{\mathscr{N}}
\newcommand{\Pt}{\mathscr{P}}
% \newcommand{\ES}{\text{\upshape\O}}
% \renewcommand{\ES}{\cancel{\mathrm{O}}}
\newcommand{\ES}{\varnothing}
\newcommand{\TopS}{\left( X, \T \right)}
\newcommand{\bcap}{\bigcap}
\newcommand{\bcup}{\bigcup}
\newcommand{\dif}{\backslash}
\newcommand{\Solution}{\noindent\textbf{Solution.}}
% \let\oldbigcup\bigcup
% \renewcommand{\bigcup}{\boldsymbol{\oldbigcup}} 

\newcommand{\ev}{\textsc{ev}}
\newcommand{\ii}{\hat{\imath}}
\newcommand{\jj}{\hat{\jmath}}
\newcommand{\kk}{\hat{k}}
\newcommand{\vv}{\mathbf{v}}
\newcommand{\va}{\mathbf{a}}
\newcommand{\vb}{\mathbf{b}}
\newcommand{\vd}{\mathbf{d}}
\newcommand{\vx}{\mathbf{x}}
\newcommand{\vy}{\mathbf{y}}
\newcommand{\vu}{\mathbf{u}}
\newcommand{\vzr}{\mathbf{0}}
\newcommand{\pp}{\boldsymbol{\cdot}}
\newcommand{\AND}{\quad\text{and}\quad}
\newcommand{\seq}{\subseteq}
\newcommand{\sep}{\supseteq}
\newcommand{\nseq}{\nsubseteq}
\newcommand{\then}{\implies}
\newcommand{\abs}[1]{\left|#1\right|}
\newcommand{\Abs}[1]{\left|\!\left|#1\right|\!\right|}
\newcommand{\prt}[1]{\left(#1\right)}
\newcommand{\End}{\hfill\(\square\)}
\newcommand{\Mod}[1]{\ \left(\mathrm{mod}\ #1\right)}
\newcommand{\floor}[1]{\left\lfloor #1 \right\rfloor}
\newcommand{\ceil}[1]{\left\lceil #1 \right\rceil}
\renewcommand\labelitemi{$\bullet$}
\newcommand{\interior}[1]{#1^\circ}
\renewcommand\labelitemi{$\bullet$}
\renewcommand{\labelitemii}{$ \circ $}
\renewcommand{\mathbb}[1]{\mathds{#1}}
\renewcommand*{\thefootnote}{\fnsymbol{footnote}}

\DeclareMathOperator{\lcm}{lcm}
\DeclareMathOperator{\adj}{adj}
\DeclareMathOperator{\sen}{sen}
\DeclareMathOperator{\proy}{proy}
\DeclareMathOperator{\tr}{tr}
\DeclareMathOperator{\spn}{span}
\DeclareMathOperator{\im}{Im}
\DeclareMathOperator{\dd}{d }
\DeclareMathOperator{\card}{card}
\DeclareMathOperator{\ddd}{d\!}
% \DeclareMathOperator{\y}{y}
\DeclareMathOperator{\rank}{rank}
\DeclareMathOperator{\nul}{nul}
\DeclareMathOperator{\Int}{int}
\DeclareMathOperator{\cla}{cla}

% Añadidos después
\newcommand{\vect}[1]{\left( {#1}_{1},\dots,{#1}_{n} \right)}

\usepackage{amsthm}
\newtheoremstyle{uptheorem} 
{0.75cm}{0.75cm}{\upshape}{}{\bfseries}{.}{ }{} 
\newtheoremstyle{sltheorem} 
{0.75cm}{0.3cm}{\slshape}{}{\bfseries}{.}{ }{} 
% \theoremstyle{sltheorem}
\theoremstyle{uptheorem}
\newtheorem{theorem}{Problema}%[chapter]
\newtheorem{definition}{Definition}%[chapter]
\newcommand{\definitionautorefname}{Definition}
\newtheorem{lemma}{Lemma}%[theorem]
\newcommand{\lemmaautorefname}{Lemma}
\newtheorem{corollary}{Corollary}%[theorem]
\newcommand{\corollaryautorefname}{Corollary}
\newtheorem{corolario}{Corolario}[theorem]
\newcommand{\corolarioautorefname}{Corolario}
\theoremstyle{uptheorem} 
\newtheorem{example}{Example}%[chapter]
  \newaliascnt{examples}{example}
  \newtheorem{examples}[examples]{Examples}
  \aliascntresetthe{examples}
  \def\exampleautorefname{Example}
  \def\examplesautorefname{Examples}

\newcommand{\notaautorefname}{Nota}

\newtheoremstyle{uptheoremv2}{0.7cm}{}{\upshape}{}{\bfseries}{.}{ }{} 
\theoremstyle{uptheoremv2} 
\newtheorem{problem}[theorem]{Problema}
\newcommand{\problemautorefname}{Problem}
\newtheoremstyle{uptheoremv3}{}{}{\upshape\small}{}{\bfseries}{.}{ }{} 
\theoremstyle{uptheoremv3}
\newtheorem{remark}{Remark}[chapter]

\numberwithin{figure}{section}
\numberwithin{table}{section}
% \numberwithin{equation}{chapter}

% \usepackage[backend=biber]{biblatex}
% \bibliography{bib.bib}
% \nocite{*}

\color{black!85}

\begin{document}
\def\thetitle{Tarea 02}
\def\fechaentrega{26 de enero de 2025}


% \pagebreak
% \vspace*{-1.5cm}
% \begin{center}
%     \includegraphics[height=0.075\textheight]{images/LogoClubMate.pdf} 
%         % \hspace{1cm}
%     %     \hfill
%     % \includegraphics[height=0.075\textheight]{images/LogoECMC.pdf}
% \end{center}

% \phantomsection\pdfbookmark{\currfilebase}{\currfilebase}

% \begin{center}
%     {\LARGE
%     Topología \\ Curso Vacacional Enero 2025\\
%     % Teaching Assistant
%     \vspace{0.5cm}
%     \textbf{\thetitle{}}}

%     % \emph{I hear, I forget;
%     % I see, I remember;
%     % I do, I understand.}
    
%     Christian Chávez\\
%     \today
% \end{center}

% % \setcounter{problem}{0}


%=================================

\begin{center}
\begin{minipage}{0.45\textwidth}
    \includegraphics[height=0.075\textheight]{images/LogoClubMate.pdf}
\end{minipage}%
\hfill
\begin{minipage}{0.5\textwidth}
    \begin{flushright}
        {\LARGE Topología} \\ Curso Vacacional -- Enero 2025 \\
        Instructor: Christian Chávez \\
    \end{flushright}
\end{minipage}

\vspace{\baselineskip}
{\huge\thetitle} \\
% \vspace{0.5cm}

\today\\
Fecha de entrega: \fechaentrega{}

\end{center}



\section*{1.\enspace Continuidad }

\begin{problem}
Demuestra que 
\begin{enumerate}[label=(\roman*)]
    \item La composición de funciones continuas es continua. 
    \item Si \(\mathcal{T}\) y \(\mathcal{U}\) son dos topologías sobre un conjunto \(X\), entonces 
        \[1_X \colon (X,\mathcal{T}) \to (X,\mathcal{U})\]
    es continua si y solamente si \(\mathcal{U}\subset \mathcal{T}\).
\end{enumerate}
\end{problem}

\begin{problem}[Equivalencias de la definición de continuidad]
Demuestra que una función \(f\colon X\to Y\)
es continua si y solo si 
\begin{enumerate}[label=(\roman*)]
    \item para todo \(A\subset Y\) ocurre que \(\text{Cl}\,  f^{-1}(A)\subset f^{-1}(\text{Cl} \, A)\)
    \item dada \(\mathcal{B}\)  una base para la topología de \(Y\), el conjunto \(f^{-1}(B)\) es abierto en \(X\) para todo \(B\in \mathcal{B}\).
\end{enumerate}
\end{problem}

\begin{problem}
Considera el mapa 

\[
f : [0, 2] \to [0, 2] \quad:\quad f(x) = 
\begin{cases} 
x & \text{si } x \in [0, 1), \\ 
3 - x & \text{si } x \in [1, 2].
\end{cases}
\]
Demuestra que \(f\) no es continua.
\end{problem}


\section*{2.\enspace Homeomorfismos}


\begin{problem}
Sea \( f: X \to Y \) una biyección continua. Son equivalentes:
\begin{enumerate}[label=(\roman*)]
    \item \( f \) es un homeomorfismo.
    \item \( f \) es abierta.
    \item \( f \) es cerrada.
\end{enumerate}


\end{problem}

\begin{problem}
 Sea \( f : X \to Y \) un homeomorfismo. Entonces, para todo \( A \subset X \),
\begin{enumerate}[label=(\roman*)]
    \item \( f(\operatorname{Cl} A) = \operatorname{Cl}(f(A)) \)
    \item \( f(\operatorname{Int} A) = \operatorname{Int}(f(A)) \)
    \item \( f(\operatorname{Fr} A) = \operatorname{Fr}(f(A)) \)
    \item \( A \) es un entorno de un punto \( x \in X \) si y solo si \( f(A) \) es un entorno del punto \( f(x) \)
\end{enumerate}
\end{problem}

\section*{3.\enspace Convergencia y continuidad en espacios métricos}

Sean \(  (X, d) \) y \( (Y, \rho) \)
espacios métricos.

\begin{problem}
Sea \( f : X \to Y \).
Demuestra que son equivalentes:
\begin{enumerate}[label=(\roman*)]
    \item \( f \) es continua sobre \( X \) (con la definición de espacios métricos)
    \item Para cualquier abierto \( V \subset Y \), el conjunto  \( f^{-1}(V) \) es abierto en \( X \)
\end{enumerate}

\end{problem}

\noindent
\textbf{Definición.}
Una función \( f : (X, d) \to (Y, \rho) \) es \textbf{lipschitz} si existe una constante \( L > 0 \) tal que 
\[
\rho(f(x), f(y)) \leq L \, d(x, y), \quad \forall x, y \in X.
\]

\begin{problem}
Demuestra que si una función es Lipschitz, entonces es continua.
\end{problem}


\end{document}

